\begin{abstract}
\setstretch{1.0}
\noindent
A machine that makes many different parts needs many different tools but can hold only a
few at a time, so running a set of jobs means stopping to change tools, and each stop
costs production time. The Job Sequencing and Tool Switching Problem (SSP) asks for the
order of jobs, and the loading at each step, that makes the total number of tool changes
smallest. The loading for a fixed order is easy; the order is NP-hard, and it is where
all of the difficulty sits.

In practice the SSP is not attacked directly: jobs are first grouped into as few
magazine-fitting batches as possible, and the batches are then sequenced.
\citet{CramaKlundert1999} showed this two-phase method to be a $b$-approximation, for
$b$ the magazine capacity. This report sharpens that into an instance-sensitive theory.
The method admits two readings the literature has not separated --- the cost of a walk
through one fixed magazine per group, and the cost once its final step re-optimises the
loading --- and the two differ on the instances most often used as examples. Within that
setting the optimum equals the cheapest walk over \emph{all} groupings, so the gap ---
measured on the best value the construction can reach, not on one implementation of it
--- is exactly the price of insisting on the fewest groups. Zero-gap classes follow for
every capacity, upper bounds for every instance, and the extremal behaviour at three
groups. In the censuses run here positive gaps are rare, every one observed equals one,
and none occurs at capacity two; the only larger-gap construction found is
disconnected, and its unboundedness remains conjectural. That regime agrees with the one
published industrial study.

On the exact side the obstacle is a single quantity. Of the compact formulations
surveyed here, the strongest relaxes to the number of tools minus the magazine capacity,
an elementary counting bound. Two methods are built to get past it: a
branch-and-Benders-cut algorithm, which settles the loading exactly with a polynomial
oracle and so never forms a weak relaxation, and position-indexed formulations, one of
which can provably exceed the counting bound.

Both were benchmarked on $1{,}421$ instances under one protocol against three compact
models reimplemented here. That baseline is not the literature at its strongest ---
Formulation~4 of \citet{Catanzaro2015} stands in for the Formulation~5 its authors
recommend --- and it closes more instances than anything built here regardless, so the
campaign yields a diagnosis and not a faster solver. Whether the counting bound equals
the optimum is the dominant predictor of what the Benders solver closes: $94\%$ of the
instances where it is tight against $40\%$ of the rest, and where it fails it has
usually already found the optimum and cannot certify it. That split is not evidence that
the loose instances are hard --- an arc-based compact model closes every one of them,
and a model whose relaxation \emph{is} the counting bound loses only thirteen points
across it. What it measures is a master carrying that bound and, at the root, little
else.

The limitation is structural rather than incidental, in this sense. What is proved: the
objective counts re-insertions produced by the global interleaving of the schedule;
apart from one constant coverage row every inequality either method carries is supported
on a single arc or a single position; and two independent configuration relaxations sit
exactly on the counting bound. That this locality is what defeats the search is the
reading those facts support, not a further theorem. A family that is not local is
derived and proved valid, with a ceiling over integral schedules several times that of
the family the solvers carry today --- but the ceiling is not reached. Measured inside
the master it leaves the root relaxation unchanged on all $37$ instances for which both
root values were obtained. The obstacle is the encoding of a global inequality, not the
absence of one.
\end{abstract}
